\subsection{Exact-v16 selected-row morphology diagnostics}
\label{app:v4p3-exact-v16-diagnostics}

The preceding kinematic atlas used a v11 SIMP/TM production because the exact
event-level source of the v4.2 observed-limit histogram was not then locally
available.  A new v16 diagnostic production now supplies the unchanged
prompt-target-constrained, $p_{\rm sum}>2.8$ baseline together with exhaustive
topology and stored-trigger partitions and four mass-versus-control
histograms.  This closes the earlier \emph{histogram-level} sample mismatch:
the 2000-, 4000-, and 8000-bin invariant-mass histograms, including their
axes, are bitwise identical to the frozen v4.2 2021 input and each contains
141,590,915 selected rows.  Unique physical-event counting remains outside
this histogram-level closure because the parser did not identify a stable
run/event key.

\begin{table}[H]
\centering
\small
\caption{Release-level continuity checks for the exact-v16 diagnostic
product.  ``Exact'' refers to histogram arrays and axes, not to an
independently frozen event-file list or unique-event identity.}
\label{tab:v4p3-v16-provenance}
\begin{tabularx}{0.96\linewidth}{P{0.31\linewidth} P{0.24\linewidth} Y}
\toprule
gate & result & qualification \\
\midrule
canonical baseline, 2000/4000/8000 bins & bitwise exact; 141,590,915 rows &
matches the frozen v4.2 spectrum at all three binnings \\
six topology categories & exact binwise partition closure &
includes the 25.6\% \texttt{topology\_other} category \\
four stored-trigger flag categories & exact binwise partition closure &
the ``neither'' category is empty; categories are exclusive stored-flag combinations \\
1-MeV mass-retention sidecar & 11 selections $\times$ 260 bins exact &
all values reproduce aggregations of the 8000-bin ROOT histograms \\
input-file provenance & 1189 unique paths; list hash recorded &
status remains \texttt{list-hash-unfrozen} \\
row semantics & selected candidate rows &
unique-event mapping and inter-row correlations are unavailable \\
\bottomrule
\end{tabularx}
\end{table}

The four controls are pair momentum sum, vertex-projection significance,
the minimum-track $|y_0|$ proxy, and the stored electron-track 3D-hit count.
The last quantity is not the N2D-hit variable used in earlier selection
studies, and vertex-projection significance is not reconstructed vertex
$z$.  No new vertex-$\chi^2$, momentum, opening-angle, or $y_0$ cut is made.
The control-axis overflows are retained in the ledger; the displayed 2D
domains omit 395 vertex-significance rows and 7368 $y_0$ rows, while the
mass axis omits 82,471 rows below 40~MeV and 256,936 above 300~MeV.

\subsubsection{Sideband-trained conditional-composition GPR}
\label{app:v4p3-composition-gpr-method}

The available histograms do not support an event-level two-dimensional
signal likelihood.  In particular, there is no signal template in the second
coordinate, no category-specific efficiency or mass-resolution model, and no
validated factorization of selected rows into independent events.  The useful
part of the true-muonium 2D-GPR construction is therefore adapted only as a
background-morphology diagnostic: a target-masked smooth surface is fit to
the mass dependence of the conditional control composition.

For mass bin $i$ and one of $J$ mutually represented control groups, let
$N_{ij}$ be the selected-row count and $N_i=\sum_jN_{ij}$.  The fitted
coordinates are the Jeffreys-regularized centered log ratios
\begin{equation}
 p_{ij}=\frac{N_{ij}+1/2}{N_i+J/2},\qquad
 y_{ij}=\log p_{ij}-\frac{1}{J}\sum_{k=1}^{J}\log p_{ik}.
 \label{eq:v4p3-clr-gpr}
\end{equation}
Independent mass-axis RBF predictors for the $y_{ij}$ coordinates are
coupled again by the inverse centered-log-ratio transform, so every predicted
composition sums to one.  This conditions on the observed mass-bin total and
tests morphology only; it neither estimates the total background under a
candidate nor assigns a signal yield.

The primary card uses 0.5-MeV mass bins over 40--160~MeV and eight contiguous
control groups whose boundaries are frozen from candidate-masked sideband
rows.  The eight stored 3D-hit bins, six exhaustive topology categories, and
three nonempty trigger categories remain discrete.  The union of the fixed
$65\pm2.5\sigma_m$ and $80\pm2.5\sigma_m$ stripes is excluded from both
kernel choice and training to protect the interpolation against reciprocal
model-mediated cross-response.  Such cross-response can arise through the
common nonlocal GP/background profile and adjacent deficit/background
structure.  It is assigned to the shared fitting procedure rather than to
physical causation, shared events, or intrinsic dependence between the
features.  For each surface the RBF length
and
standardized nugget are chosen from
\begin{equation}
 \ell_m\in\{2,4,8,12,20,35\}\ \mathrm{MeV},\qquad
 \eta\in\{10^{-4},10^{-3},10^{-2},5\times10^{-2}\},
\end{equation}
using the median centered-log-ratio residual in five same-rule held-out fake
gaps at 50, 100, 115, 130, and 145~MeV.  None of those scores is converted to
a $p$-value or significance.  Deterministic checks repeat the procedure at
0.125, 0.25, 0.5, and 1.0~MeV mass binning, with 6, 8, 10, and 12 continuous
control groups, and with single-target rather than union masks.  No
pseudoexperiment, expected band, or coverage claim is introduced in v4.3.

\begin{figure}[H]
\centering
\graphicorplaceholder{0.98\linewidth}{final_limit_projection_figs/v4p3_2021_10pct_exact_v16_diagnostics_20260810/partition_composition_gpr.pdf}
\caption{Exact-v16 topology and stored-trigger compositions aggregated over
the 65 and 80~MeV target stripes.  Black bars are observed selected-row
fractions and orange bars are the jointly masked, sideband-trained
conditional-composition GPR predictions.  The exact partition closures are a
data-integrity result.  The displayed interpolation is a correlated
histogram-level morphology diagnostic, not a category signal fit.}
\label{fig:v4p3-partition-composition}
\end{figure}

\subsubsection{Fixed-window controls and mass-surface residuals}
\label{app:v4p3-control-results}

The exact parser windows contain 16,361,818 rows at 65~MeV and 20,984,092 rows
at 80~MeV.  Subtracting the combined adjacent sidebands after the prescribed
$5/3$ width scaling gives $+1,039,881$ and $+557,880$ rows, respectively.
Those large values are spectrum-curvature diagnostics: 65~MeV lies on the
rising shoulder of a distribution whose broad maximum is near 77.5~MeV.
They are not fitted candidate yields.  After normalizing the central and
combined-sideband control distributions separately and applying the primary
control grouping, the total-variation distances span only 0.12--1.02\% across
the four saved variables.  Lower and
upper sidebands move in opposite directions for several controls, confirming
that a one-number sideband subtraction is mass sculpted.

\begin{figure}[H]
\centering
\graphicorplaceholder{0.98\linewidth}{final_limit_projection_figs/v4p3_2021_10pct_exact_v16_diagnostics_20260810/central_sideband_controls.pdf}
\caption{Normalized exact-v16 control distributions in each central stripe
and its lower, upper, and combined adjacent sidebands.  Each curve is
normalized independently to expose shape rather than the rapidly changing
mass yield.  The visible lower-to-upper evolution motivates the
mass-conditioned GPR; none of these differences is a fitted signal
component.}
\label{fig:v4p3-central-sideband-controls}
\end{figure}

\begin{figure}[H]
\centering
\graphicorplaceholder{0.98\linewidth}{final_limit_projection_figs/v4p3_2021_10pct_exact_v16_diagnostics_20260810/conditional_residual_atlas.pdf}
\caption{Observed minus GPR-predicted conditional composition in the 65 and
80~MeV target stripes for the four mass--control surfaces and the topology and
trigger partitions.  Colors are percentage-point composition residuals, not
Poisson pulls.  All panels reuse the same selected rows and must not be
combined as independent tests.}
\label{fig:v4p3-conditional-residual-atlas}
\end{figure}

The parent v4.2 deterministic inpainting supplies a useful scale for the
dilution of any local component.  In its complete analysis-bin windows,
$(N_{\rm data}-N_{\rm GP\ mean})/N_{\rm data}$ is 0.2605\% at 65~MeV and
0.1460\% at 80~MeV.  These are conditional mean-fill differences, not signal
purities.  The parent counterfactual changes each complete window total by
less than 0.3\%, which illustrates the severe dilution scale but does not
identify the remaining rows as background.  Sub-percent composition changes
therefore cannot by themselves distinguish a dilute prompt signal from
ordinary mass dependence or reconstruction effects.

\begin{figure}[H]
\centering
\graphicorplaceholder{0.98\linewidth}{final_limit_projection_figs/v4p3_2021_10pct_exact_v16_diagnostics_20260810/effect_size_robustness_summary.pdf}
\caption{Conditional-composition effect sizes and deterministic closure
context at 65 and 80~MeV.  The upper panels show the primary maximum
composition shift and its ratio to the median fake-gap CLR residual; the
ratio is not a $p$-value.  In the lower panels, filled markers show the
primary union-masked card and line segments span the mass-binning variants.
The ranges are model-robustness envelopes, not confidence intervals.  Parent
v4.2 mean-fill fractions are included only as dilution references.}
\label{fig:v4p3-effect-size-summary}
\end{figure}

The 65~MeV stripe has modest, repeatable changes in some momentum-sum,
topology, and high-$|y_0|$ fractions, while the trigger composition is nearly
unchanged.  The detailed vertex-significance and 3D-hit-bin directions are
less stable under reasonable binning and grouping choices.  At 80~MeV the
effects are generally smaller and their leading bins or signs change more
often across the deterministic variations.  Because each comparison reuses
the same selected rows and conditions on the mass-bin total, differences
across variables or regions are correlated descriptors rather than separate
tests.  In particular, the joint mask does not convert shared-model
cross-response into a statement about the physical relation between the
65 and 80~MeV features.

Thus the exact-v16 product removes the earlier histogram-level sample
mismatch and provides a controlled view of saved-variable morphology.  The
histograms alone do not furnish an event-level signal template, fitted yield,
or calibrated local or global significance.  Event-identified run/time
splits, category-specific signal simulation and response, and
correlation-preserving inference would be required for those quantities.
